\documentclass[10pt,a4paper]{amsart}
\usepackage{amsmath,amssymb,amsthm,graphicx}
\usepackage{verbatim}
\newenvironment{reference}{\comment}{\endcomment}
%\newenvironment{reference}{}{}
\newenvironment{slogan}{\comment}{\endcomment}
\newenvironment{history}{\comment}{\endcomment}

% For commutative diagrams we use Xy-pic
\usepackage[all]{xy}

% We use 2cell for 2-commutative diagrams.
\xyoption{2cell}
% Theorem environments.
%
\theoremstyle{plain}
\newtheorem{theorem}[subsection]{Theorem}
\newtheorem{proposition}[subsection]{Proposition}
\newtheorem{lemma}[subsection]{Lemma}

\theoremstyle{definition}
\newtheorem{definition}[subsection]{Definition}
\newtheorem{example}[subsection]{Example}
\newtheorem{exercise}[subsection]{Exercise}
\newtheorem{situation}[subsection]{Situation}

\theoremstyle{remark}
\newtheorem{remark}[subsection]{Remark}
\newtheorem{remarks}[subsection]{Remarks}

\numberwithin{equation}{subsection}

% Macros
%
\def\lim{\mathop{\mathrm{lim}}\nolimits}
\def\colim{\mathop{\mathrm{colim}}\nolimits}
\def\Spec{\mathop{\mathrm{Spec}}}
\def\Hom{\mathop{\mathrm{Hom}}\nolimits}
\def\Ext{\mathop{\mathrm{Ext}}\nolimits}
\def\SheafHom{\mathop{\mathcal{H}\!\mathit{om}}\nolimits}
\def\SheafExt{\mathop{\mathcal{E}\!\mathit{xt}}\nolimits}
\def\Sch{\mathit{Sch}}
\def\Mor{\mathop{\mathrm{Mor}}\nolimits}
\def\Ob{\mathop{\mathrm{Ob}}\nolimits}
\def\Sh{\mathop{\mathit{Sh}}\nolimits}
\def\NL{\mathop{N\!L}\nolimits}
\def\CH{\mathop{\mathrm{CH}}\nolimits}
\def\proetale{{pro\text{-}\acute{e}tale}}
\def\etale{{\acute{e}tale}}
\def\QCoh{\mathit{QCoh}}
\def\Ker{\mathop{\mathrm{Ker}}}
\def\Im{\mathop{\mathrm{Im}}}
\def\Coker{\mathop{\mathrm{Coker}}}
\def\Coim{\mathop{\mathrm{Coim}}}

% Boxtimes
%
\DeclareMathSymbol{\boxtimes}{\mathbin}{AMSa}{"02}

%
% Macros for moduli stacks/spaces
%
\def\QCohstack{\mathcal{QC}\!\mathit{oh}}
\def\Cohstack{\mathcal{C}\!\mathit{oh}}
\def\Spacesstack{\mathcal{S}\!\mathit{paces}}
\def\Quotfunctor{\mathrm{Quot}}
\def\Hilbfunctor{\mathrm{Hilb}}
\def\Curvesstack{\mathcal{C}\!\mathit{urves}}
\def\Polarizedstack{\mathcal{P}\!\mathit{olarized}}
\def\Complexesstack{\mathcal{C}\!\mathit{omplexes}}
% \Pic is the operator that assigns to X its picard group, usage \Pic(X)
% \Picardstack_{X/B} denotes the Picard stack of X over B
% \Picardfunctor_{X/B} denotes the Picard functor of X over B
\def\Pic{\mathop{\mathrm{Pic}}\nolimits}
\def\Picardstack{\mathcal{P}\!\mathit{ic}}
\def\Picardfunctor{\mathrm{Pic}}
\def\Deformationcategory{\mathcal{D}\!\mathit{ef}}
\usepackage{cleveref}
\crefname{originallemma}{original lemma}{original lemmas}
\Crefname{originallemma}{Original lemma}{Original lemmas}
\makeatletter
\newlabel{schemes-lemma-widetilde-pullback}{{01I9}{source}{}{sourcepointer-1}{}}
\newlabel{schemes-lemma-widetilde-pullback@cref}{{[originallemma][0][]01I9}{[1][0][]source}{}{}{}}
\makeatother
\title{Every quasi-coherent module is a directed colimit of kappa-generated quasi-coherent submodules}
\author{The Stacks Project Authors}
\date{stacks-project:properties.tex:local-snapshot-20261012; GNU FDL 1.2 or later}
\begin{document}
\maketitle
\noindent\textit{Editorial note.} All delivered mathematical text is the author's own, copied from the pinned Stacks Project chapter file properties.tex; no statement or proof step was rewritten, reordered or removed.\par\medskip
\noindent\textit{Editorial note.} The counted proof is the author's complete argument: the choice of the covering data with a cardinal bounding all index sets, the refinement operator on families of subsets together with its iteration to a stable family and the glueing to a kappa-generated subsheaf, and the directedness and exhaustiveness of the resulting system.\par\medskip
\noindent\textit{Editorial note.} The label lemma-colimit-kappa names the counted statement, and the proof environment immediately following it in the source is delivered verbatim as its proof.\par\medskip
\noindent\textit{Editorial note.} The counted statement is carried by the author's own lemma environment in the chapter source.\par\medskip
\noindent\textit{Editorial note.} Pre-open gate: python operations/agent-authoring/tools/precard\_gate.py --source properties.tex --label lemma-colimit-kappa --start 3771 --end 3848 --local-id 28039 returned RESULT=OK (L28039 free, occupancy FREE, public interval clear). Candidate taken from sources/stacks-properties/candidates-halfA.json.\par\medskip
\noindent\textit{Editorial note.} The statement and proof are copied from the pinned chapter file; the cited inputs are kept as external pointers with their canonical Stacks tags.\par\medskip
\noindent\textit{Allowed standard prerequisites (editorial).}\begin{itemize}
\item Base change of sections of a quasi-coherent module along an inclusion of affine opens. Schemes, the lemma identifying the sections of the restriction of a quasi-coherent module with the tensor product of the sections over the smaller affine open. Supplies the identification M\_i tensor over A\_i with A\_ijk = M\_ijk = M\_j tensor over A\_j with A\_ijk used to define the refinement map.
\item Description of quasi-coherent modules on a scheme by their values on an affine open covering. Schemes, the section describing quasi-coherent modules through the values on the members of an affine open covering and the glueing data. Supplies the construction of a quasi-coherent subsheaf from a compatible family of submodules.
\end{itemize}\medskip
\ifdefined\hypertarget\hypertarget{sourcepointer-1}{}\fi
\noindent\textit{Source reference.} Original lemma 01I9: sources/stacks-chow/tags.txt, source lines 1970--1970. Prerequisite pointer invoked at proof line(s) 3793; it names a standard fact of a sibling chapter that the counted argument uses, and every step of the counted proof built on it is retained verbatim in the delivered body.\par\medskip
\begin{lemma}
\label{lemma-colimit-kappa}
Let $X$ be a scheme. There exists a cardinal $\kappa$ such that
every quasi-coherent module $\mathcal{F}$ is the directed colimit
of its quasi-coherent $\kappa$-generated submodules.
\end{lemma}

\begin{proof}
Choose an affine open covering $X = \bigcup_{i \in I} U_i$. For each pair
$i, j$ choose an affine open covering
$U_i \cap U_j = \bigcup_{k \in I_{ij}} U_{ijk}$.
Write $U_i = \Spec(A_i)$ and $U_{ijk} = \Spec(A_{ijk})$.
Let $\kappa$ be any infinite cardinal $\geq$ than the cardinality
of any of the sets $I$, $I_{ij}$.

\medskip\noindent
Let $\mathcal{F}$ be a quasi-coherent sheaf. Set $M_i = \mathcal{F}(U_i)$
and $M_{ijk} = \mathcal{F}(U_{ijk})$. Note that
$$
M_i \otimes_{A_i} A_{ijk} = M_{ijk} = M_j \otimes_{A_j} A_{ijk}.
$$
see
Schemes, Lemma \ref{schemes-lemma-widetilde-pullback}.
Using the axiom of choice we choose a map
$$
(i, j, k, m) \mapsto S(i, j, k, m)
$$
which associates to every $i, j \in I$, $k \in I_{ij}$ and $m \in M_i$
a finite subset $S(i, j, k, m) \subset M_j$ such that we have
$$
m \otimes 1 = \sum\nolimits_{m' \in S(i, j, k, m)} m' \otimes a_{m'}
$$
in $M_{ijk}$ for some $a_{m'} \in A_{ijk}$. Moreover, let's agree
that $S(i, i, k, m) = \{m\}$ for all $i, j = i, k, m$ as above.
Fix such a map.

\medskip\noindent
Given a family $\mathcal{S} = (S_i)_{i \in I}$ of subsets
$S_i \subset M_i$ of cardinality at most $\kappa$ we set
$\mathcal{S}' = (S'_i)$ where
$$
S'_j = \bigcup\nolimits_{(i, k, m)\text{ such that }m \in S_i}
S(i, j, k, m)
$$
Note that $S_i \subset S'_i$. Note that $S'_i$ has cardinality at most
$\kappa$ because it is a union over a set of cardinality at most $\kappa$
of finite sets. Set $\mathcal{S}^{(0)} = \mathcal{S}$,
$\mathcal{S}^{(1)} = \mathcal{S}'$ and by induction
$\mathcal{S}^{(n + 1)} = (\mathcal{S}^{(n)})'$. Then set
$\mathcal{S}^{(\infty)} = \bigcup_{n \geq 0} \mathcal{S}^{(n)}$.
Writing $\mathcal{S}^{(\infty)} = (S^{(\infty)}_i)$ we see that
for any element $m \in S^{(\infty)}_i$ the image of $m$ in
$M_{ijk}$ can be written as a finite sum $\sum m' \otimes a_{m'}$
with $m' \in S_j^{(\infty)}$. In this way we see that setting
$$
N_i = A_i\text{-submodule of }M_i\text{ generated by }S^{(\infty)}_i
$$
we have
$$
N_i \otimes_{A_i} A_{ijk} = N_j \otimes_{A_j} A_{ijk}.
$$
as submodules of $M_{ijk}$. Thus there exists a quasi-coherent subsheaf
$\mathcal{G} \subset \mathcal{F}$ with $\mathcal{G}(U_i) = N_i$.
Moreover, by construction the sheaf $\mathcal{G}$ is $\kappa$-generated.

\medskip\noindent
Let $\{\mathcal{G}_t\}_{t \in T}$ be the set of $\kappa$-generated
quasi-coherent subsheaves. If $t, t' \in T$ then
$\mathcal{G}_t + \mathcal{G}_{t'}$ is also a $\kappa$-generated
quasi-coherent subsheaf as it is the image of the map
$\mathcal{G}_t \oplus \mathcal{G}_{t'} \to \mathcal{F}$.
Hence the system (ordered by inclusion) is directed.
The arguments above show that every section of $\mathcal{F}$ over $U_i$
is in one of the $\mathcal{G}_t$ (because we can start with $\mathcal{S}$
such that the given section is an element of $S_i$). Hence
$\colim_t \mathcal{G}_t \to \mathcal{F}$ is both injective and surjective
as desired.
\end{proof}

\end{document}
